\documentclass[12pt,a4paper]{article}
\usepackage[utf8]{inputenc}
\usepackage[T1]{fontenc}
\usepackage{amsmath,amssymb,amsthm}
\usepackage{geometry}
\usepackage{hyperref}
\usepackage{booktabs}

\geometry{margin=1in}

\newtheorem{theorem}{Theorem}
\newtheorem{proposition}{Proposition}
\newtheorem{conjecture}{Conjecture}
\newtheorem{assumption}{Assumption}
\theoremstyle{remark}
\newtheorem{remark}{Remark}

\title{\textbf{Gaussian temporal profile of gravitational decoherence\\[6pt]
from the Einstein-Langevin equation}}

\author{Will Regelmann\thanks{Corresponding author.} \quad Claude (Anthropic)\thanks{AI
assistant; contributed to mathematical development and analysis.}}

\date{February 2026}

\begin{document}
\maketitle

\begin{abstract}
Several experimental programs are approaching the sensitivity needed to test gravitational
decoherence at the mesoscopic scale. The primary theoretical target, the
Di\'osi--Penrose timescale $\tau_{DP} = \hbar/E_\Delta$, is derived either by postulating
white noise in the Newtonian gravitational potential (Di\'osi) or by invoking the
ill-definedness of the time-translation operator for superposed geometries (Penrose).
Neither route uses the Einstein equation. We ask what the Einstein-Langevin
equation---the established stochastic extension of semiclassical gravity---actually
predicts for the decoherence temporal profile of a massive body in spatial superposition.

We derive two results and state one conjecture. (1)~The decoherence timescale is
$\tau_{\text{coh}} = (4\sqrt{2}/5)\,\hbar/E_\Delta \approx 1.13\,\tau_{DP}$, determined by
the same gravitational self-energy $E_\Delta$ as the Di\'osi--Penrose prediction. (2)~The
temporal profile is Gaussian, not exponential---a direct consequence of the noise kernel
for stationary matter being time-independent in the Newtonian limit. Both results are
derived at sketch rigor: they treat the stress-energy fluctuations of the superposition as
Gaussian noise, an assumption built into the Einstein-Langevin framework that we make
explicit, since the underlying stress-energy distribution of a two-branch superposition is
bimodal. (3)~We conjecture that the profile is material-dependent: rigid crystalline
bodies give Gaussian decoherence, while bodies whose acoustic modes approach the
gravitational coherence frequency $\omega_g = E_\Delta/\hbar$ would transition toward
exponential. If the crossover regime is experimentally accessible, comparing profiles
across materials of identical mass and geometry would provide a discriminant absent from
existing models. The precise additional postulate that Di\'osi's model makes beyond the
Einstein-Langevin equation is identified: a white-noise temporal correlator for the
gravitational potential fluctuations.
\end{abstract}

%======================================================================
\section{Introduction}
%======================================================================

Experiments using levitated nanoparticles and matter-wave interferometry are approaching
the sensitivity needed to test gravitational decoherence at mesoscopic
masses~\cite{bose,maqro}. The principal theoretical benchmark is the Di\'osi--Penrose
timescale:
\begin{equation}
\tau_{DP} = \frac{\hbar}{E_\Delta}, \qquad
E_\Delta \equiv \frac{G}{2}\iint \frac{\Delta\rho(\mathbf{x})\,\Delta\rho(\mathbf{x}')}{|\mathbf{x}-\mathbf{x}'|}\,d^3x\,d^3x'
\label{eq:tau_dp}
\end{equation}
where $\Delta\rho = \rho_L - \rho_R$ is the mass-density difference between the two
superposed configurations. Di\'osi~\cite{diosi} derives this by postulating a white-noise
stochastic process in the Newtonian potential with a specific noise amplitude. Penrose
\cite{penrose} reaches the same timescale from the ill-definedness of the time-translation
Killing vector for superposed spacetime geometries. Di\'osi's model predicts exponential decoherence,
$|\rho_{LR}(t)| \propto e^{-t/\tau_{DP}}$; Penrose's argument identifies the same
timescale $\tau_{DP}$ but does not specify a temporal profile.

Neither derivation begins from the Einstein equation or its stochastic extension. The
Einstein-Langevin equation~\cite{hu_verdaguer} is the established framework for including
quantum stress-energy fluctuations in semiclassical gravity; it has been used to derive the
spectrum of cosmological perturbations without quantizing the metric~\cite{roura_verdaguer}
and to establish a rigorous theory of noise in curved spacetime~\cite{phillips_hu}. The
question we address is: what does this framework predict for the temporal structure of
gravitational decoherence of a massive body in spatial superposition?

This question sits within an existing field-theoretic literature on gravitational
decoherence, surveyed by Bassi, Gro\ss ardt and Ulbricht~\cite{bassi_review}.
Anastopoulos and Hu~\cite{anastopoulos_hu} derive a gravitational-decoherence master
equation from quantized linear perturbations interacting with matter, obtaining
decoherence in the energy basis, and caution against treating the stress-energy variance
of a quantum state as a classical noise source; Blencowe~\cite{blencowe} reaches a
master equation of similar structure from gravity as an effective-field-theory
environment. The present paper does not supersede those first-principles treatments. Its
contribution is narrower: to take the Einstein-Langevin equation as given---including its
defining prescription that stress-energy fluctuations enter as Gaussian noise, which we
state as an explicit assumption (Assumption~\ref{ass:gaussian}) rather than a derived
fact---and to work out what that framework, so understood, predicts for the temporal
profile. The Gaussian-profile result should be read as a sharp consequence of the
Einstein-Langevin postulates, exposing exactly where Di\'osi's model adds a further
postulate, not as a field-theoretic derivation of decoherence from first principles.

The deterministic semiclassical Einstein equation alone cannot answer this question. For a
body in superposition $|\psi\rangle = (|L\rangle + |R\rangle)/\sqrt{2}$, the
stress-energy expectation value averages over branches,
$\langle\hat{T}_{\mu\nu}\rangle = \frac{1}{2}(T^{(L)}_{\mu\nu} + T^{(R)}_{\mu\nu})$,
carrying no information about the off-diagonal coherence $\rho_{LR}$. The resulting
geometry is blind to the superposition structure and cannot decohere it. The stochastic
extension is necessary.

We derive two results and state one conjecture. First, the decoherence timescale from the
Einstein-Langevin equation is:
\begin{equation}
\tau_{\text{coh}} = \frac{4\sqrt{2}}{5}\,\frac{\hbar}{E_\Delta} \approx 1.13\,\tau_{DP}
\label{eq:tau_coh_preview}
\end{equation}
within 13\% of the DP prediction and determined by the same physical quantity $E_\Delta$.
Second, the temporal profile is Gaussian:
\begin{equation}
|\rho_{LR}(t)| = |\rho_{LR}(0)|\,e^{-t^2/\tau_{\text{coh}}^2}
\label{eq:gaussian_preview}
\end{equation}
This arises because the noise kernel for stationary matter is time-independent in the
Newtonian limit, producing static rather than white noise. Both results are
\textbf{(Sketch)}: every algebraic step is given, but the chain rests on the Gaussian-noise
assumption of the Einstein-Langevin framework (Assumption~\ref{ass:gaussian}), which is
not derived. Third, we conjecture that the profile is material-dependent: for rigid
crystalline bodies the noise is static and the profile is Gaussian; for bodies with
acoustic modes near the gravitational coherence frequency $\omega_g = E_\Delta/\hbar$, a
finite noise correlation time would develop and the profile would transition toward
exponential (Conjecture~\ref{conj:material}). If an experimentally accessible platform
reaches that crossover regime---an open question---a comparison of decoherence profiles
across materials of identical mass and geometry could distinguish whether the noise is
fundamental (Di\'osi) or derived from quantum stress-energy (this work).

Section~\ref{sec:framework} recalls the Einstein-Langevin equation.
Section~\ref{sec:noise_kernel} evaluates the noise kernel for a rigid body in
superposition. Section~\ref{sec:gaussian} derives the Gaussian profile and timescale.
Section~\ref{sec:dp_comparison} identifies the precise additional postulate in Di\'osi's
model. Section~\ref{sec:predictions} states the experimental predictions.
Section~\ref{sec:discussion} discusses implications and limitations.

%======================================================================
\section{The Einstein-Langevin framework}
\label{sec:framework}
%======================================================================

The Einstein-Langevin equation is the leading-order stochastic correction to the
semiclassical Einstein equation, accounting for quantum fluctuations of the matter
stress-energy~\cite{hu_verdaguer}:
\begin{equation}
G_{\mu\nu}[g] = 8\pi G\!\left(\langle\hat{T}_{\mu\nu}\rangle_g + \xi_{\mu\nu}\right)
\label{eq:einstein_langevin}
\end{equation}
Here $G_{\mu\nu}[g]$ is the Einstein tensor, $\langle\hat{T}_{\mu\nu}\rangle_g$ is the
renormalized expectation value of the stress-energy tensor in the quantum state on the
background geometry $g$, and $\xi_{\mu\nu}$ is a real Gaussian stochastic tensor field
representing quantum stress-energy fluctuations.

The stochastic source $\xi_{\mu\nu}$ has vanishing mean,
$\langle\xi_{\mu\nu}(x)\rangle_s = 0$, and two-point function given by the noise kernel:
\begin{equation}
\langle\xi_{\mu\nu}(x)\,\xi_{\alpha\beta}(x')\rangle_s
= N_{\mu\nu\alpha\beta}(x,x')
\equiv \tfrac{1}{2}\langle\{\delta\hat{T}_{\mu\nu}(x),\,\delta\hat{T}_{\alpha\beta}(x')\}\rangle
\label{eq:noise_kernel_def}
\end{equation}
where $\delta\hat{T}_{\mu\nu} \equiv \hat{T}_{\mu\nu} - \langle\hat{T}_{\mu\nu}\rangle$
is the stress-energy fluctuation operator and $\langle\cdot\rangle$ denotes the quantum
expectation value in the physical state.

\begin{assumption}[Gaussian noise]
\label{ass:gaussian}
The stochastic source $\xi_{\mu\nu}$ is a \emph{Gaussian} random field, fully specified
by its vanishing mean and the two-point function~\eqref{eq:noise_kernel_def}.
\end{assumption}

\noindent
This is the defining prescription of the Einstein-Langevin framework~\cite{hu_verdaguer},
not a derived property, and for the states considered here it is a genuine physical
assumption: for a two-branch superposition the actual distribution of $\hat{T}_{00}$ is
bimodal (the mass is in one place or the other), which is maximally non-Gaussian. The
Einstein-Langevin prescription matches only the symmetrized two-point function of
$\delta\hat{T}_{\mu\nu}$ and discards all higher cumulants. The decoherence
formula~\eqref{eq:decoherence_integral} used below is exact only for Gaussian noise.
Every result downstream of Assumption~\ref{ass:gaussian} therefore carries at most
\textbf{(Sketch)} rigor as a statement about gravitational decoherence in nature, however
exact it is as a statement about the Einstein-Langevin framework. Field-theoretic
treatments that do not make this assumption reach qualitatively different
conclusions~\cite{anastopoulos_hu,blencowe}.

In the Newtonian (non-relativistic, weak-field)
limit, the linearized metric perturbation satisfies the stochastic Poisson equation:
\begin{equation}
\nabla^2 \delta\Phi(\mathbf{x},t) = 4\pi G\,\xi_{00}(\mathbf{x},t)/c^2
\label{eq:poisson_stochastic}
\end{equation}
with retardation corrections of order $v^2/c^2$.

%======================================================================
\section{Noise kernel for a superposed rigid body}
\label{sec:noise_kernel}
%======================================================================

\subsection{Evaluation}

Consider a rigid body of mass $m$ placed in a center-of-mass superposition
$|\psi\rangle = (|L\rangle + |R\rangle)/\sqrt{2}$, with branch densities
$\rho_L(\mathbf{x}) = \rho(\mathbf{x}-\mathbf{x}_L)$ and
$\rho_R(\mathbf{x}) = \rho(\mathbf{x}-\mathbf{x}_R)$, separated by
$d = |\mathbf{x}_L - \mathbf{x}_R|$. For non-relativistic matter, $T_{00}$ is the energy
density: $\hat{T}_{00}(\mathbf{x}) = c^2\rho(\mathbf{x}-\hat{\mathbf{X}}_{\text{CM}})$
(this convention is used consistently with the stochastic Poisson
equation~\eqref{eq:poisson_stochastic}), giving:
\begin{align}
\langle\hat{T}_{00}(\mathbf{x})\rangle &=
\tfrac{c^2}{2}\!\left[\rho_L(\mathbf{x}) + \rho_R(\mathbf{x})\right] \\
\langle\hat{T}_{00}(\mathbf{x})\hat{T}_{00}(\mathbf{x}')\rangle &=
\tfrac{c^4}{2}\!\left[\rho_L(\mathbf{x})\rho_L(\mathbf{x}')
+ \rho_R(\mathbf{x})\rho_R(\mathbf{x}')\right]
\end{align}
The cross terms $\rho_L(\mathbf{x})\rho_R(\mathbf{x}')$ and
$\rho_R(\mathbf{x})\rho_L(\mathbf{x}')$ vanish because the branches are orthogonal,
$\langle L|R\rangle = 0$: the operator
$\hat{T}_{00}(\mathbf{x})\hat{T}_{00}(\mathbf{x}')$ connects $|L\rangle$ to $|L\rangle$
and $|R\rangle$ to $|R\rangle$ (each branch has definite center-of-mass position), so only
the diagonal matrix elements survive.
Subtracting, the noise kernel~\eqref{eq:noise_kernel_def} evaluates to:
\begin{equation}
N_{0000}(\mathbf{x},\mathbf{x}') = \frac{c^4}{4}\,\Delta\rho(\mathbf{x})\,\Delta\rho(\mathbf{x}')
\label{eq:N0000}
\end{equation}
where $\Delta\rho \equiv \rho_L - \rho_R$. The noise kernel factorizes as the outer
product of the branch-difference density with itself. This evaluation is
\textbf{(Rigorous)} given the non-relativistic stress-energy operator and the branch
orthogonality stated above; what is \emph{not} rigorous is the subsequent use of this
two-point function as the complete specification of a Gaussian noise source
(Assumption~\ref{ass:gaussian}).

\subsection{Physical interpretation}

The noise is not a quantum vacuum fluctuation of matter fields. It is the statistical
variance of the mass distribution arising from the superposition: the mass is in two
places, and this uncertainty generates metric fluctuations.

\begin{remark}[Vacuum fluctuations negligible; order of magnitude only]
\label{rem:vacuum}
Quantum vacuum fluctuations of $\hat{T}_{00}$ within a single branch are suppressed
relative to eq.~\eqref{eq:N0000} by factors of order
$(\ell_P/\sigma)^2 \sim G\hbar/(c^3\sigma^2)$~\cite{phillips_hu}, where $\sigma$ is the
body size. For a micron-scale diamond ($\sigma \approx 0.9\,\mu\text{m}$, cf.\
Table~\ref{tab:platforms}) this ratio is $\sim 3\times10^{-58}$: entirely negligible.
This is an order-of-magnitude estimate, not a derived bound.
\end{remark}

\subsection{Time-independence in the Newtonian limit}

For a stationary superposition---one in which wavepacket spreading is negligible on the
decoherence timescale---the mass distributions $\rho_L$ and $\rho_R$ are
time-independent. Consequently the noise kernel is static:
\begin{equation}
N_{0000}(\mathbf{x},t;\,\mathbf{x}',t') = \frac{c^4}{4}\,\Delta\rho(\mathbf{x})\,\Delta\rho(\mathbf{x}')
\label{eq:N0000_static}
\end{equation}
Retardation corrections are $\mathcal{O}(v^2/c^2)$ and negligible. The static character
of eq.~\eqref{eq:N0000_static} is the key structural feature that distinguishes the
Einstein-Langevin noise from the white noise postulated by Di\'osi.

The stationary approximation holds when the wavepacket spreading time $md^2/\hbar$ exceeds
the decoherence timescale. For a microdiamond ($m \sim 10^{-14}$~kg, $d \sim 250\,\mu$m),
this timescale is $md^2/\hbar \approx 6\times10^{12}$~s versus $\tau_{\text{coh}} \approx
13$~ms (Table~\ref{tab:platforms})---a margin of $\sim 5\times10^{14}$.

%======================================================================
\section{Gaussian decoherence timescale}
\label{sec:gaussian}
%======================================================================

\subsection{Decoherence integral}

The stochastic Newtonian potential at branch location $\mathbf{x}_{L,R}$ is:
\begin{equation}
\delta\Phi(\mathbf{x}) = -\frac{G}{c^2}\int \frac{\xi_{00}(\mathbf{x}')}{|\mathbf{x}-\mathbf{x}'|}\,d^3x'
\end{equation}
The branch-differential stochastic force $\delta V_\Delta(t) \equiv m[\delta\Phi(\mathbf{x}_L,t) - \delta\Phi(\mathbf{x}_R,t)]$ dephases the off-diagonal element of the density matrix. Averaging over stochastic realizations~\cite{hu_verdaguer}:
\begin{equation}
|\rho_{LR}(t)| = |\rho_{LR}(0)|\exp\!\left(-\frac{1}{2\hbar^2}\int_0^t\!\!\int_0^t
C_V(t'-t'')\,dt'\,dt''\right)
\label{eq:decoherence_integral}
\end{equation}
where $C_V(\tau) \equiv \langle\delta V_\Delta(t)\,\delta V_\Delta(t+\tau)\rangle_s$.

\subsection{Static noise gives Gaussian suppression}

From eqs.~\eqref{eq:N0000_static} and~\eqref{eq:poisson_stochastic}, the correlator
inherits time-independence: $C_V(\tau) = C_V(0)$ for all $\tau$. The double integral in
eq.~\eqref{eq:decoherence_integral} then evaluates to $C_V(0)\,t^2$, giving:
\begin{equation}
|\rho_{LR}(t)| = |\rho_{LR}(0)|\,\exp\!\left(-\frac{C_V(0)}{2\hbar^2}\,t^2\right)
\label{eq:gaussian_suppression}
\end{equation}
This is the direct mathematical consequence of static noise. In the Markovian limit, a
finite noise correlation time $\tau_c$ gives $\int\int C_V dt'dt'' \approx 2C_V(0)\tau_c
t$, which through the $1/(2\hbar^2)$ prefactor in eq.~\eqref{eq:decoherence_integral}
yields exponential decoherence at rate $\Gamma = C_V(0)\tau_c/\hbar^2$. In the
static limit ($\tau_c\to\infty$ for stationary matter), the exponential rate vanishes and
the profile is Gaussian with timescale $\tau_{\text{coh}} = \hbar\sqrt{2}/\sqrt{C_V(0)}$.

\begin{remark}[Ensemble dephasing versus single-system decoherence]
\label{rem:dephasing}
Equation~\eqref{eq:gaussian_suppression} is an average over realizations of a
\emph{static} Gaussian random variable. In any single realization the branch-differential
potential $\delta V_\Delta$ is a constant, and the accumulated phase is deterministic and
in principle recoverable (e.g.\ by echo protocols); the Gaussian decay describes
dephasing of an ensemble. Di\'osi's white noise, by contrast, produces irreversible
decoherence realization by realization. Whether the ensemble-averaged Gaussian suppression
constitutes genuine single-system decoherence---and what, physically, the ensemble of
realizations of the superposition's own stress-energy variance refers to---is an open
interpretational question of the Einstein-Langevin framework that we flag here and do not
resolve.
\end{remark}

\subsection{Evaluation of $C_V(0)$ for a uniform sphere}

Propagating the factored noise kernel $N_{0000}\propto\Delta\rho\,\Delta\rho'$
through the Poisson Green's function, the stochastic-potential two-point function
factorizes as $\langle\delta\Phi(\mathbf{x})\,\delta\Phi(\mathbf{y})\rangle_s \propto U(\mathbf{x})\,U(\mathbf{y})$
where $U(\mathbf{x})\equiv \int\Delta\rho(\mathbf{z})/|\mathbf{x}-\mathbf{z}|\,d^3z$.
Consequently $\langle[\delta\Phi(\mathbf{x}_L)-\delta\Phi(\mathbf{x}_R)]^2\rangle_s
\propto [U(\mathbf{x}_L)-U(\mathbf{x}_R)]^2$, and assembling the overall prefactor:
\begin{equation}
C_V(0) = \frac{G^2 m^2}{4}\!\left[U(\mathbf{x}_L) - U(\mathbf{x}_R)\right]^2, \qquad
U(\mathbf{x}) \equiv \int\frac{\Delta\rho(\mathbf{z})}{|\mathbf{x}-\mathbf{z}|}\,d^3z
\label{eq:CV0}
\end{equation}
For a uniform sphere of radius $R$ with $d > 2R$ (non-overlapping branches), Newton's
shell theorem gives\footnote{The potential
$\int\rho(\mathbf{z})/|\mathbf{x}_L - \mathbf{z}|\,d^3z$ of a uniform sphere of mass $m$
and radius $R$, evaluated at its own center, equals $3m/(2R)$: in spherical coordinates,
$\int_0^R \rho_0 \cdot 4\pi r^2/r\,dr = 2\pi\rho_0 R^2 = 3m/(2R)$.  Since
$\Delta\rho = \rho_L - \rho_R$ and the branches do not overlap,
$U(\mathbf{x}_L) = 3m/(2R) - m/d$ (self-potential of the $L$-sphere at its center, minus
the external potential of the $R$-sphere at distance $d$). By symmetry
$U(\mathbf{x}_R) = -3m/(2R) + m/d$, giving the stated result.}:
\begin{equation}
U(\mathbf{x}_L) - U(\mathbf{x}_R) = m\!\left(\frac{3}{R} - \frac{2}{d}\right)
\end{equation}
Substituting into eq.~\eqref{eq:CV0} and recognizing
$E_\Delta = G(6m^2/5R - m^2/d)$~(eq.~\ref{eq:tau_dp}):
\begin{equation}
C_V(0) = \frac{25}{16}\,E_\Delta^2\left(1 + \mathcal{O}(R/d)\right)
\label{eq:CV0_sphere}
\end{equation}
\noindent
The leading correction\footnote{Writing $\varepsilon = R/d$ and
$f(\varepsilon) \equiv C_V(0)/E_\Delta^2 = (3-2\varepsilon)^2/[4(6/5-\varepsilon)^2]$,
one finds $f(0) = 25/16$ and $(\mathrm{d}\ln f/\mathrm{d}\varepsilon)|_{\varepsilon=0} = 1/3$, giving
$C_V(0) = \tfrac{25}{16}E_\Delta^2(1 + R/(3d) + \mathcal{O}(R^2/d^2))$ and
$\tau_{\text{coh}} = (4\sqrt{2}/5)(\hbar/E_\Delta)(1 - R/(6d) + \mathcal{O}(R^2/d^2))$.}
is positive in $C_V(0)$ and negative in $\tau_{\text{coh}}$.
The $G^2$ factor in $C_V(0)$ arises from the two-step chain: stress-energy noise converts
to metric perturbation (one factor of $G$), which converts to potential (a second). For
the overlapping regime ($d < 2R$), $E_\Delta$ must be computed via the Fourier-space
integral:
\begin{equation}
E_\Delta = \frac{2Gm^2}{\pi}\int_0^\infty|\tilde{F}(kR)|^2\!\left(1-\frac{\sin kd}{kd}\right)dk, \quad
\tilde{F}(kR) = \frac{3[\sin kR - kR\cos kR]}{(kR)^3}
\label{eq:E_delta_fourier}
\end{equation}

\subsection{Coherence timescale}

\begin{equation}
\tau_{\text{coh}} = \frac{\hbar\sqrt{2}}{\sqrt{C_V(0)}}
= \frac{4\sqrt{2}}{5}\,\frac{\hbar}{E_\Delta}\left(1 + \mathcal{O}(R/d)\right)
\approx 1.13\,\tau_{DP} \quad (d \gg R)
\label{eq:tau_coh}
\end{equation}
The $G^2$ in $C_V(0)$ yields $G^{-1}$ scaling in $\tau_{\text{coh}}$ via the square root,
matching the $G$-scaling of $\tau_{DP}$.

\textbf{Rigor status of results (1) and (2).} Equations~\eqref{eq:gaussian_suppression}
and~\eqref{eq:tau_coh} are \textbf{(Sketch)}. Every algebraic step from the noise
kernel~\eqref{eq:N0000} to the timescale is given above, and within the Einstein-Langevin
framework the chain is complete; but the chain rests on
Assumption~\ref{ass:gaussian} (Gaussian noise, not derived, and maximally violated by the
actual bimodal stress-energy distribution of the state), on the decoherence
formula~\eqref{eq:decoherence_integral}, which is exact only for Gaussian noise and is
imported from the stochastic-gravity literature~\cite{hu_verdaguer} rather than derived
here, and on the interpretational caveat of Remark~\ref{rem:dephasing}. Numerical values
for three experimental platforms are given in Table~\ref{tab:platforms}.

\begin{table}[h]
\centering
\begin{tabular}{lcccccc}
\toprule
Platform & $m$ & $R$ & $d$ & $\tau_{DP}$ & $\tau_{\text{coh}}$ & Profile \\
\midrule
Microdiamond (BMV)~\cite{bose} & $10^{-14}$~kg & 880~nm & 250~$\mu$m & 11.6~ms & 13.1~ms & Gaussian \\
Levitated silica nanosphere & $5\times10^{-18}$~kg & 82~nm & 500~nm & 5000~s & 5500~s & Gaussian \\
MAQRO test mass~\cite{maqro} & $1.7\times10^{-17}$~kg & 123~nm & 100~nm & 2880~s & 2850~s & Gaussian \\
\bottomrule
\end{tabular}
\caption{Coherence timescales for three experimental platforms. Radii follow from the
masses and bulk densities (diamond $3500$~kg/m$^3$; fused silica $2200$~kg/m$^3$). The
MAQRO entry is in the overlapping-branch regime ($d/R \approx 0.8$): $E_\Delta$ (and hence
$\tau_{DP}$) is computed numerically from eq.~\eqref{eq:E_delta_fourier}, while $C_V(0)$
is computed from eq.~\eqref{eq:CV0} using the interior potential of a uniform sphere,
which gives $U(\mathbf{x}_L)-U(\mathbf{x}_R) = md^2/R^3$ for $d \le R$ and hence
$\tau_{\text{coh}} = 2\sqrt{2}\,\hbar R^3/(Gm^2d^2)$. In this regime the prefactor
$4\sqrt{2}/5$ in eq.~\eqref{eq:tau_coh} does not apply ($\tau_{\text{coh}} \approx
0.99\,\tau_{DP}$ here). All three platforms satisfy
$\omega_{\text{phonon}}/\omega_g \gg 1$ and are in the rigid-body Gaussian regime.}
\label{tab:platforms}
\end{table}

%======================================================================
\section{Relation to the Di\'osi--Penrose prediction}
\label{sec:dp_comparison}
%======================================================================

\subsection{Same timescale}

Both the Einstein-Langevin derivation and the Di\'osi--Penrose prediction give a timescale
set by $E_\Delta$. This is not coincidental: the bilinear structure $\langle\Delta\rho,\Delta\rho'\rangle_{1/r}$
of the noise kernel selects $E_\Delta$ as the natural energy scale in the Newtonian limit. The specific prefactor
$4\sqrt{2}/5 \approx 1.13$ in eq.~\eqref{eq:tau_coh} comes from the uniform-sphere
geometry and could differ for other mass distributions.

\subsection{Different temporal correlator: the extra postulate}

The profiles differ because the noise correlators differ. The Einstein-Langevin equation
gives a static correlator (eq.~\ref{eq:N0000_static}), which propagates to:
\begin{equation}
C_V^{\text{EL}}(\tau) = C_V(0) \quad \forall\,\tau
\label{eq:corr_el}
\end{equation}
Di\'osi's model postulates a white-noise correlator:
\begin{equation}
C_V^{\text{D}}(\tau) = 2C_V(0)\,\tau_c\,\delta(\tau)
\label{eq:corr_diosi}
\end{equation}
with $\tau_c$ chosen so that $\Gamma = C_V(0)\tau_c/\hbar^2 = 1/\tau_{DP}$ reproduces
the DP rate. The spatial structure and equal-time amplitude $C_V(0)$ are identical in both
cases; the difference is entirely in the temporal correlator. Di\'osi's extra postulate,
beyond what the Einstein-Langevin equation specifies, is the Markovian (white-noise)
character of the temporal correlations.

A summary comparison with the Kafri--Taylor--Milburn model~\cite{kafri} is given in
Table~\ref{tab:models}.

\begin{table}[h]
\centering
\begin{tabular}{lllll}
\toprule
Model & Timescale & $G$-scaling & Profile & Noise temporal correlator \\
\midrule
This work & $\hbar/E_\Delta$ & $1/G$ & Gaussian & Static (derived) \\
Di\'osi--Penrose & $\hbar/E_\Delta$ & $1/G$ & Exponential & White (postulated) \\
Kafri--Taylor--Milburn & $\sim \omega d^3/(G m)$ & $1/G$ & Exponential & White (postulated) \\
\bottomrule
\end{tabular}
\caption{Comparison of gravitational decoherence models. The Kafri--Taylor--Milburn
model~\cite{kafri} applies to two masses $m$ in harmonic traps of frequency $\omega$ at
separation $d$: its decoherence rate $\Lambda = K/(4m\omega)$, with gravitational spring
constant $K = 2Gm^2/d^3$, is of the order of the gravitationally induced normal-mode
splitting, giving a timescale $\sim m\omega/K \sim \omega d^3/(Gm)$. This work and
Di\'osi--Penrose share the same timescale scaling; the discriminating prediction is the
temporal profile.}
\label{tab:models}
\end{table}

%======================================================================
\section{Experimental predictions}
\label{sec:predictions}
%======================================================================

\subsection{Profile discrimination}

The maximum fractional difference between the predicted profiles occurs near
$t \approx \tau_{DP}/2$. Using $|\rho_{LR}^{\text{Gauss}}| = e^{-t^2/\tau_{\text{coh}}^2}$
with $\tau_{\text{coh}} = 1.13\,\tau_{DP}$, and
$|\rho_{LR}^{\text{exp}}| = e^{-t/\tau_{DP}}$:
\begin{equation}
\max_t\frac{|\rho_{LR}^{\text{Gauss}}| - |\rho_{LR}^{\text{exp}}|}{|\rho_{LR}(0)|}
\approx 0.22
\label{eq:delta_V}
\end{equation}
This maximum occurs at $t \approx 0.48\,\tau_{DP} \approx 0.43\,\tau_{\text{coh}}$,
where the Gaussian has decayed by $\sim18\%$ and the exponential by $\sim40\%$.
In the limiting approximation $\tau_{\text{coh}} \approx \tau_{DP} \equiv \tau$,
the difference at $t = \tau/2$ simplifies to $e^{-1/4} - e^{-1/2} \approx 0.17$.
To resolve the actual 0.22 difference at $3\sigma$ requires per-time-point visibility
uncertainty $\sigma_V < 0.07$. For shot-noise-limited measurements with fringe contrast
$C \approx 0.5$ this demands $N_{\text{run}} \gtrsim 1000$ shots per time point. Across
eight time points spanning $t \in [0.1\tau_{DP}, 3\tau_{DP}]$, the total is
$\sim 8{,}000$ experimental shots.

The profile discrimination requires that environmental decoherence $\Gamma_{\text{env}}$
be characterized to within a factor of $\sim 2$ (so that $\delta\Gamma_{\text{env}}\cdot
\tau_{\text{coh}} \lesssim 0.1$). This is achievable by comparing data at different masses
or between materials with the same environmental coupling but different $E_\Delta$.

\subsection{Material-dependent profile transition (conjectural)}
\label{sec:material}

The static-noise result~\eqref{eq:N0000_static} holds when the body's lowest acoustic
phonon frequency $\omega_{\text{phonon}}$ far exceeds the gravitational coherence frequency
$\omega_g = E_\Delta/\hbar$. For the micron-scale diamond of Table~\ref{tab:platforms},
$\omega_{\text{phonon}} \sim c_s/L \sim 10^{10}$~rad/s while
$\omega_g \approx 90$~rad/s: the ratio
$\omega_{\text{phonon}}/\omega_g \sim 10^8$ confirms that the rigid-body Gaussian result
applies to all current experimental targets.

\begin{conjecture}[Material-dependent crossover]
\label{conj:material}
When $\omega_{\text{phonon}} \sim \omega_g$, internal acoustic dynamics introduce a finite
correlation time $\tau_c \sim \omega_{\text{phonon}}^{-1}$ into the noise kernel, and the
decoherence exponent crosses over:
\begin{equation}
\int_0^t\!\!\int_0^t C_V(t'-t'')\,dt'dt'' \approx
\begin{cases}
C_V(0)\,t^2 & t \ll \tau_c \\
2\,C_V(0)\,\tau_c\,t & t \gg \tau_c
\end{cases}
\label{eq:crossover}
\end{equation}
so that bodies of identical mass and geometry but lower rigidity drift from a Gaussian
toward an exponential profile.
\end{conjecture}

\noindent
This is a \textbf{(Conjecture)}: the identification $\tau_c \sim
\omega_{\text{phonon}}^{-1}$ is asserted, not derived---a phonon-spectrum model of the
noise kernel is an open problem (Section~\ref{sec:discussion}). Moreover, whether the
crossover regime is experimentally accessible at all is open. For a mesoscopic body of
size $L$, $\omega_{\text{phonon}} \sim c_s/L$; even soft materials ($c_s \sim
10^2$--$10^3$~m/s for aerogels and polymers) give $\omega_{\text{phonon}} \sim
10^8$--$10^{10}$~rad/s at $L \sim 1\,\mu$m, still many orders of magnitude above
$\omega_g \lesssim 10^2$~rad/s for the masses of Table~\ref{tab:platforms}. Reaching
$\omega_{\text{phonon}} \sim \omega_g$ requires much larger bodies or much larger
$E_\Delta$; quantifying whether any realistic platform reaches it is the open problem
noted above. If the regime is accessible, the discriminant is sharp: the Di\'osi model and
the Penrose collapse model predict exponential decoherence regardless of material, so
observing profile differences across materials at fixed mass and geometry would directly
probe whether the noise originates in quantum stress-energy (this work) or is a
fundamental background process.

\subsection{BMV entanglement}

In the Bose--Marletto--Vedral experiment~\cite{bose,marletto_vedral}, gravitational
entanglement generation between two superposed masses would signal quantum gravity. In the
Einstein-Langevin framework, gravity mediates via a classical stochastic channel. For a
bipartite product state $\rho_1\otimes\rho_2$, the metric perturbation factorizes as
$h_{\mu\nu} = h^{(1)}_{\mu\nu} + h^{(2)}_{\mu\nu}$, yielding a c-number interaction
potential that acts as a classical common-cause channel. Because entanglement cannot
be generated by local operations on a separable state coupled only through a classical
stochastic field~\cite{kafri,oppenheim,marletto_classical}, the prediction is no
entanglement at the quantum gravity rate; observation of such entanglement would falsify
the framework. This prediction is \textbf{(Sketch)}: the factorization argument above is
given only in outline, and its force is conditional on the cited no-go results for
entanglement generation through classical channels, which are used here as established
inputs rather than re-derived.

%======================================================================
\section{Discussion}
\label{sec:discussion}
%======================================================================

\textbf{What the derivation assumes.} The results rest on five assumptions within
the Hu-Verdaguer stochastic gravity program: (i)~the Einstein-Langevin equation describes
metric fluctuations sourced by quantum stress-energy; (ii)~the Newtonian limit applies;
(iii)~the superposition is stationary on the decoherence timescale; (iv)~leading-order
perturbation theory (backreaction of decoherence on the noise kernel is neglected);
(v)~the stress-energy fluctuations enter as \emph{Gaussian} noise
(Assumption~\ref{ass:gaussian})---the framework's defining prescription, not a derived
fact, and the reason results (1) and (2) are labeled Sketch.
Backreaction becomes significant when the compactness parameter $Gm/(c^2 R)$ approaches
unity, i.e., when the body's gravitational radius is comparable to its size. For all
laboratory masses this ratio is negligible: $Gm/(c^2 R) \sim 10^{-35}$ for a
microdiamond. Perturbation theory breaks down only for compact objects with
$R \sim Gm/c^2$, far outside the experimental regime considered here.
No physics beyond the semiclassical Einstein equation is introduced.

\textbf{State-dependent noise.} The noise kernel~\eqref{eq:N0000} vanishes when
$\Delta\rho = 0$ (no superposition). The metric fluctuations that decohere the
superposition are generated by the superposition itself, not by a background stochastic
process. This is physically consistent and distinguishes the present framework
from Di\'osi's model, where the noise is a fundamental property of spacetime. It also has
an observational consequence: the underground constraints on the Di\'osi--Penrose model
from spontaneous radiation emission~\cite{donadi} probe a universal background noise that
acts on every charged particle at all times; a state-dependent noise that vanishes in the
absence of superposition predicts no such radiation from unsuperposed bulk matter and is
not constrained by those experiments.

\textbf{Limitations.} The Gaussian profile is a leading-order prediction; retardation
corrections are $\mathcal{O}(v^2/c^2)$. The material-dependent crossover
(Conjecture~\ref{conj:material}) requires a more detailed model of the phonon-spectrum
contribution to the noise kernel, and a quantitative assessment of whether any
experimentally realistic platform reaches the crossover regime; both are open problems.
The ensemble-dephasing caveat of Remark~\ref{rem:dephasing} is unresolved. The BMV null
prediction is necessary but not sufficient; the discriminating positive prediction is the
temporal profile and, conditionally on Conjecture~\ref{conj:material}, its material
dependence.

\textbf{Companion paper.} The semiclassical Einstein equation underlying the
Einstein-Langevin formalism is itself a fixed-point condition: geometry and quantum state
must be mutually self-consistent. The existence of self-consistent solutions---established
exactly for the Starobinsky inflation model and conditionally in general via the Schauder
fixed-point theorem---is treated in a companion paper~\cite{fixed_point}.

%======================================================================
\begin{thebibliography}{99}

\bibitem{bose} S.~Bose \textit{et al.}, Spin entanglement witness for quantum gravity,
\textit{Phys.\ Rev.\ Lett.}\ \textbf{119}, 240401 (2017).

\bibitem{maqro} R.~Kaltenbaek \textit{et al.}, Macroscopic quantum resonators (MAQRO),
\textit{Exp.\ Astron.}\ \textbf{34}, 123 (2012).

\bibitem{diosi} L.~Di\'osi, A universal master equation for the gravitational violation
of quantum mechanics, \textit{Phys.\ Lett.\ A}\ \textbf{120}, 377 (1987).

\bibitem{penrose} R.~Penrose, On gravity's role in quantum state reduction,
\textit{Gen.\ Rel.\ Grav.}\ \textbf{28}, 581 (1996).

\bibitem{hu_verdaguer} B.~L.~Hu and E.~Verdaguer, Stochastic gravity: Theory and
applications, \textit{Living Rev.\ Relativ.}\ \textbf{11}, 3 (2008).

\bibitem{roura_verdaguer} A.~Roura and E.~Verdaguer, Cosmological perturbations from
stochastic gravity, \textit{Phys.\ Rev.\ D}\ \textbf{78}, 064010 (2008).

\bibitem{phillips_hu} N.~G.~Phillips and B.~L.~Hu, Noise kernel in stochastic gravity and
stress energy bitensor, \textit{Phys.\ Rev.\ D}\ \textbf{63}, 104001 (2001).

\bibitem{anastopoulos_hu} C.~Anastopoulos and B.~L.~Hu, A master equation for
gravitational decoherence: probing the textures of spacetime,
\textit{Class.\ Quantum Grav.}\ \textbf{30}, 165007 (2013).

\bibitem{blencowe} M.~P.~Blencowe, Effective field theory approach to gravitationally
induced decoherence, \textit{Phys.\ Rev.\ Lett.}\ \textbf{111}, 021302 (2013).

\bibitem{bassi_review} A.~Bassi, A.~Gro\ss ardt, and H.~Ulbricht, Gravitational
decoherence, \textit{Class.\ Quantum Grav.}\ \textbf{34}, 193002 (2017).

\bibitem{donadi} S.~Donadi, K.~Piscicchia, C.~Curceanu, L.~Di\'osi, M.~Laubenstein, and
A.~Bassi, Underground test of gravity-related wave function collapse,
\textit{Nat.\ Phys.}\ \textbf{17}, 74 (2021).

\bibitem{marletto_vedral} C.~Marletto and V.~Vedral, Gravitationally induced entanglement
between two massive particles is sufficient evidence of quantum effects in gravity,
\textit{Phys.\ Rev.\ Lett.}\ \textbf{119}, 240402 (2017).

\bibitem{kafri} D.~Kafri, J.~M.~Taylor, and G.~J.~Milburn, A classical channel model for
gravitational decoherence, \textit{New J.\ Phys.}\ \textbf{16}, 065020 (2014).

\bibitem{oppenheim} J.~Oppenheim, A postquantum theory of classical gravity?,
\textit{Phys.\ Rev.\ X}\ \textbf{13}, 041040 (2023).

\bibitem{marletto_classical} C.~Marletto, J.~Oppenheim, V.~Vedral, and E.~Wilson,
Classical gravity cannot mediate entanglement,
preprint arXiv:2511.07348 (2025).

\bibitem{fixed_point} W.~Regelmann and Claude (Anthropic), Fixed-point existence of
self-consistent semiclassical gravity (2026). Companion paper.

\end{thebibliography}

\end{document}
